\documentclass[11pt,a4paper]{article}
\title{hdb++}
\author{Lorenzo Pivetta}
\date{\today}

\usepackage{svn-multi}
\svnRegisterAuthor{pivetta}{Lorenzo Pivetta}
\svnid{$Id$}

\usepackage{amsmath}
\usepackage{amsthm}
\usepackage{amsfonts}
\usepackage{amssymb}
\usepackage{color}
\usepackage[dvips]{graphicx}
\usepackage{fancyhdr}
\usepackage{float}
\usepackage{hangcaption}
\usepackage{latexsym}
% \usepackage[none]{hyphenat}

\usepackage{longtable}

% \usepackage{draftwatermark}
% \SetWatermarkText{Draft}
%% \SetWatermarkFontSize{4cm}
% \SetWatermarkScale{6}

\topmargin			0mm
\textheight 		220mm
\textwidth			170mm
\evensidemargin 0mm
\oddsidemargin  0mm
\headheight 		0mm


\pagestyle{fancyplain}

\newcommand{\data}{19.09.2014}
\newcommand{\titolo}{\bf HDB++}

\renewcommand{\familydefault}{\sfdefault}
\renewcommand{\sectionmark}[1]{\markright{\thesection \ \ #1}}

% \lhead[\fancyplain{}{\rm\thepage}]{\fancyplain{}{\it\rightmark}}
% \rhead[\fancyplain{}{\it\leftmark}]{\fancyplain{}{\rm\thepage}}

\renewcommand{\headheight}{1.5cm}
\setlength{\headwidth}{\textwidth}
\fancyhead[L]{ % left
	\includegraphics[height=0.53in]{logo_elettra_medium.ps}
}
\fancyhead[C]{ % center
	\includegraphics[height=0.53in]{logo_tango.ps}
}
\fancyhead[R]{ % right
  \includegraphics[height=0.53in]{logo_esrf.ps}
}
\pagestyle{fancy}

\lfoot[\fancyplain{}{\footnotesize}]{\fancyplain{}{\textit{
	HDB++ meeting - \data{} - Rev. \svnrev
}}}
\rfoot[\fancyplain{}{\footnotesize}]{\fancyplain{}{\textit{
	\thepage
}}}
\cfoot{}

\begin{document}

% \maketitle

\begin{table}[htb]
	\center
	\begin{tabular}{|p{2.5cm}|p{12.5cm}|}
		\hline
		\bf{Meeting} & \bf{HDB++} \\
		\hline
		\bf{Date} & \data \\
		\hline
		\bf{Time/Place} & Teleconf \\
		\hline		
		\bf{Contact} & \parbox[][1.5cm][c]{10cm}{
			Lorenzo Pivetta \\
			lorenzo.pivetta@elettra.eu \\
			++39 040 3758361
		} \\
		\hline
		\bf{Partecipants} &  \parbox[][4.0cm][c]{10cm}{
			R. Bourtembourg - RB - ESRF \\
			JM.Chaize - JMC - ESRF \\
			L.Pivetta - LP - ELETTRA \\
			C.Scafuri - CS - ELETTRA \\
			G.Scalamera - GSC - ELETTRA \\
			P. Verdier - PV - ESRF \\
			L. Zambon - LZ - ELETTRA \\
		} \\
		\hline
		\bf{Agenda} & HDB++ design guidelines revision \\
		\hline
	\end{tabular}
\end{table}

\tableofcontents
% \listoffigures
\listoftables

\newpage

\section{Requirements/developments for the next future}

\begin{longtable}{|c|p{6.0cm}|p{6.0cm}|p{1.0cm}|p{1.4cm}|}
	\hline
	\bf{\#} & \bf{Request} & \bf{To do} & \bf{Who} & \bf{When} \\
	\hline
	\endhead
	1.1 &
	In case of error, store the error description &
	Add varchar column to the data\_type tables and related code
	to EventSubscriber &
	Elettra &
	Available \\
	\hline
	1.2 &
	Store quality factor in addition to data value &
	Add column to the data\_type tables and related code
	to EventSubscriber &
	Elettra &
	Nov 2014 \\
	\hline
	1.3 &
	Make some attribute configuration parameters available
	in the historical database &
	Add a table to store \emph{display-unit}, \emph{format-string}
	and \emph{label}. The table will also contain att\_conf\_id and
	the three timestamps.
	Add a fifo and producer/consumer threads to subscribe
	to attribute-configuration-change events into EventSubscriber &
	Elettra &
	Dec 2014 \\
	\hline
	1.4 &
	Multiple TANGO host support: Tango 8 &
	The correct FQDN has to be specified or just the
	domain/family/member can be specified letting the
	ConfigurationManager select the right TANGO\_HOST/port &
	Elettra &
	Available \\
	\hline
	1.5 &
	Support for write-only attribute &
	Bugfix &
	Elettra &
	Available \\
	\hline
	1.6 &
	Additional indexing on att\_name varchar in att\_conf table
	for faster search &
	Modify att\_conf table &
	Elettra &
	Dec 2014 \\
	\hline
	1.7 &
	Historical database data extraction library &
	Core library written in C++; C wrapper to be used
	also with php and python. Define the list of interface methods
	and send to ESRF ASAP &
	Elettra &
	Prototype and docs available \\
	\hline
	1.8 &
	GUI for Configuration Manager &
	java based GUI tool &
	ESRF &
	Available\\
	\hline
	1.9 &
	Configure polling, event parameters and alarm threshold
	programmatically for AttributeNokNumber in EventSubscriber
	to enable alarms &
	&
	Elettra &
	Dec 2014 \\
	\hline
	1.10 &
	Do not count stopped attributes as faulty &
	&
	Elettra &
	Dec 2014 \\
	\hline
	1.11 &
	Provide per-attribute errors list &
	Is last error/exception description sufficient? Or do we need
	to store the error stack/history? &
	t.b.d. &
	\\
	\hline
	1.12 &
	Provide max value for AttributePendingNumber &
	Implement AttributeMaxPendingNumber attribute &
	Elettra &
	Dec 2014 \\
	\hline
	1.13 &
	Provide number-of-events counter per attribute; reset by ResetStatistics;
	store the timestamp of the last call to ResetStatistics in an attribute
	(seconds since EPOCH) &
	Already available with the Attribute Status command; should this be:
	\begin{itemize}
		\item a command with FQDN as input value
		\item an attribute vector type containing all the counters
	\end{itemize}
	AttributeEvenNumberList &
	Elettra &
	Available\\
	\hline
	\caption{Requirements/developments for the next future}
	\label{tab:next-future}
\end{longtable}

\newpage

\section{Requirements/developments for the not-so-near future}

\begin{longtable}{|c|p{6.0cm}|p{6.0cm}|p{1.0cm}|p{1.4cm}|}
	\hline
	\bf{\#} & \bf{Request} & \bf{To do} & \bf{Who} & \bf{When} \\
	\hline
	\endhead
	2.1 &
	Multiple TANGO host support: Tango 9 & 
	Use the DbGetCSDbServerList command to verify the FQDN.
	Then we still face the problem how to manage the aliases; there are
	two possibilities:
	\begin{itemize}
		\item the ConfigurationManager can replace the alias with the
					real hostname; not the desired behaviour
		\item the DatabaseDS has to be modified in order to return also
					the alias name as a valid FQDN; the events will come from
					the proper FQDN (i.e. alias)
	\end{itemize} &
	Elettra, ESRF &
	To be clarified. Is it necessary? \\
	\hline
	2.2 &
	Monitor on periodic event based archiving &
	Provide timeout on periodic event subscriptions; can be
	implemented as an alarm timer in the configuration change
	event thread &
	Elettra &
	March 2015 \\
	\hline
	2.3 &
	Array support. Currently arrays are stored as fixed length
	strings into varchar; this may lead to loss of precision.
	Alternate approaches:
	\begin{itemize}
		\item store into native data type using an additional index to
		    	reconstruct the array
		\item store in binary form
		\item use native array support
		\item do nothing and rely to timestamp
	\end{itemize} &
	Test postgres native array support &
	Elettra, ESRF &
	Available with Cassandra. Evaluate Postgres and all... \\
	\hline
	2.4 &
	Device locking mechanism to avoid concurrent setups &
	Collaborative client/server locking mechanism &
	ESRF, Elettra &
	Available in CM GUI \\
	\hline
	2.5 &
	java-based extraction library &
	Pure java implementation of the same API as the C/C++ extraction
	library; depends on the C API &
	ESRF &
	Prototype available \\
	\hline
	2.6 &
	Addititonal database table to store data statistics for scalar
	values &
	New table with start\_time, stop\_time, min, max, average;
	additional statistics may be useful; write a tango device server,
	based on the extraction library, to calculate and store the
	statistics; depends on the extraction library &
	ESRF, Elettra &
	March 2015 \\
	\hline
	2.7 &
	Extraction library data format support &
	Attribute history data format, json data format, both?!? &
	Elettra, ESRF &
	t.b.d. \\
	\hline
	2.8 &
	NoSQL databases &
	NoSQL databases evaluation &
	ESRF &
	Prototype available \\
	\hline
	2.9 &
	Data stream management system (DSMS) &
	InfluxDB evaluation &
	Elettra &
	March 2015 \\
	\hline
	2.10 &
	Store encrypted passwords in the device/class property &
	&
	Elettra &
	March 2015 \\
	\hline
	\caption{Requirements/developments for the not-so-near future}
	\label{tab:next-but-future}
\end{longtable}

\newpage

\section{Pending stuff}

\begin{longtable}{|c|p{6.0cm}|p{6.0cm}|p{1.0cm}|p{1.4cm}|}
	\hline
	\bf{\#} & \bf{Request} & \bf{To do} & \bf{Who} & \bf{When} \\
	\hline
	\endhead
	3.1 &
	Polling thread often/always late on last attribute of the poller
	configuration &
	&
	&
	\\
	\hline
	3.1 &
	dev\_enum: dedicated data table to store value and label; both value
	and configuration change event subscription &
	&
	&
	\\
	\hline
	3.3 &
	Archiving support for device server commands (TANGO core releated) &
	&
	&
	\\
	\hline
	\caption{Pending stuff}
	\label{tab:pending}
\end{longtable}

\section{Best practices}

\begin{longtable}{|c|p{6.0cm}|p{6.0cm}|p{1.0cm}|p{1.4cm}|}
	\hline
	\bf{\#} & \bf{Request} & \bf{To do} & \bf{Who} & \bf{When} \\
	\hline
	\endhead
	4.1 &
	Run/Shutdown support - 1st &
	Separate attributes that have to be archived all the time from
	those that aren't (e.g. during shutdown) using different dedicated 
	EventSubscriber instances; this allows for straightforward
	start/stop archiving on per-instance basis (e.g. from a higher layer
	Tango device implementing some sort of machine shifts calendar) &
	&
	\\
	\hline
	4.2 &
	Run/Shutdown support  - 2nd&
	Introduce per-attribute flag to specify stopped-at-shutdown; an
	additional command StopForShutdown will stop all flagged attributes.
	To maintain the whole configuration into the Tango database the
	flag will be specified comma separated in the AttributeList property. &
	Elettra &
	\\
	\hline
	4.3 &
	Old data filtering &
	Add "managed" flag to the att\_conf table, updated by the AttributeAdd and
	AttributeRemove methods &
	t.b.d. &
	\\
	\hline
	\caption{Best practices}
	\label{tab:best-practice}
\end{longtable}

\end{document}








